

The successful chronic implantation of next-gen BCIs is constrained by extreme power efficiency requirements due to patient safety concerns. This has historically been met through highly inflexible, monolithic ASIC hardware. HALO BCIs diverge from the norm by using various function-specific PEs and the RVV extension to introduce significant flexibility without compromising critical thermal limits. However, a highly optimized software compilation pipeline is required to make HALO and its clinical flexibility and capabilities accessible to the broader neuroscientific community. This project directly addresses this need by engineering an end-to-end toolchain that merges a TypeScript-based source-to-source compiler with a hardware-specific vectorization library, enabling the automatic transformation of high-level MATLAB algorithms into power-optimized RISC-V binaries.

Validation and benchmarking demonstrate that the vectorized matrix library successfully runs on Spike and yields near-consistent performance improvements. Analysis of the data has successfully isolated performance improvements driven by updated cross-compiler toolchains, while simultaneously identifying initial test-harness anomalies. The ts-traversal semantic translation layer has also been benchmarked and timed with successful results across various essential operations.

% Future research can take multiple directions from here. First, there are unvectorized kernels flagged in the RISCV-Matrix README that require further work. Second, the hand-off to the BCI lab at the Yale School of Medicine will move the pipeline from simulator-only validation onto the physical HALO prototype under neuroscientist-written full workloads, where the 15 mW power envelope and $1 ^\circ$C thermal limit will provide a real-world check on the cycle-count argument made here.

\section{Future Work}

The work presented here opens several concrete directions for future development, organized by component.

\paragraph{Continued vectorization.}
The RISCV-Matrix README flags dozens of kernels that remain unvectorized. The highest value vectorizations are likely trigonometric functions like \texttt{sinM} and the complex side of \texttt{argM} would provide significant speedups for BCI workloads due to the DSP-heavy nature of neural decoding. As discussed earlier, the vectorization of \texttt{carg} would speed-up the heavy Octave workloads from Table~\ref{tab:total-runtime}. However, this is likely a difficult endeavor as the lack of RVV trigonometric intrinsics requires that polynomial approximations to be built or pulled from  a library like SLEEF \cite{sleef}. Outside these options, \texttt{popvarM} can be vectorized easily following the implementations of \texttt{varM,stdM,popstdM}. Finally, 1D non-double types can be vectorized for \texttt{varM}.

\paragraph{Integer overflow in \texttt{mpowerM}.}
As noted in the C-Octave Suite Residuals discussion, \texttt{power\_matrix} produces results exceeding $\texttt{INT\_MAX} = 2^{31} - 1$ for \texttt{INTEGER} inputs. The Octave suite correctly computes these large values while the C suite gets stuck at \texttt{INT\_MIN/INT\_MAX}. This warrants further investigation into \texttt{mpowerM} and \texttt{floorM}, both of which get called in the test suite.

\paragraph{LMUL audit.}
Several vectorized functions were flagged in the original README as having LMUL settings that ``need checking.'' Using an LMUL that exceeds the number of live vector variables causes implicit register spills, negating the throughput benefit. An LMUL audit pass would identify and correct sub-optimal groupings.

\paragraph{Transpiler improvements.} 
Support for MATLAB's \texttt{classdef} OOP is incomplete and would be a highest value change in order to enable supporting BCI workloads through EEGLAB: the \texttt{getClasses} pass discovers properties and methods but emits \texttt{initVar()} calls inside C \texttt{struct} definitions, which is illegal~C. Separating struct declaration (header) from constructor initialization (function) would unblock the \texttt{class\_test} suite.

\paragraph{HALO deployment.}
The hand-off to the BCI lab at the Yale School of Medicine will enable physical HALO testing under neuroscientist-written full workloads, where the 15~mW power envelope and $1\,^{\circ}$C thermal limit will provide a real-world check on the cycle-count conclusions drawn here. Cache effects, TLB pressure, and memory bandwidth saturation---none of which are modeled by \texttt{rdinstret}-based benchmarking on Spike---will then become the dominant performance factors to measure and analyze.

